\etitle{Running example.}
The question asks for the mean number of identity-theft reports among
metropolitan areas whose 2023 population exceeds one million.  The agent
declares 17 CSV files---one national identity-theft table and 16 state-level
identity-theft tables---but its final program $P^*$ actually opens the
national table and 16 state-level \emph{fraud-and-other} tables.  It uses the
latter only to construct an admissible set of metropolitan-area names,
back-computes population as reports divided by reports per 100,000 times
$100{,}000$, retains 14 rows above the population threshold, and averages
their report counts.  Host execution prints
\texttt{\{"main-task": 10724.1429\}}, and the agent states the same value
$\hat a=10{,}724.1429$ (the benchmark reference answer is $12{,}964.8727$).
Under restricted replay, only the canonicalized declared files are visible.
The 16 fraud-and-other paths therefore fail the program's existence checks;
$P^*$ continues, forms an empty match, and terminates normally with
\texttt{\{"main-task": null\}}.  The replay has neither an execution
exception nor an empty stdout, but it supplies no scalar comparable with
$\hat a$.  The audit therefore withholds a replay certificate and records a
missing/non-comparable-output witness together with the authorization failure
$F_{\mathrm{host}}\nsubseteq F_{\mathrm{res}}$.

\begin{verbatim}
valid_names = union(read_msa_names(16 undeclared fraud CSVs))
rows = national_identity[norm_msa in valid_names]
rows["population"] = rows["reports"] / rows["rate_per_100k"] * 100000
answer = mean(rows[rows.population > 1_000_000]["reports"])
\end{verbatim}
